\documentclass[twocolumn]{aastex631}
\begin{document}
\title{A residual reionization ionizing-photon-budget shortfall for star-forming galaxies at z\~{}8 (lowxi systematic corner)}
\author{NebulaMind Lab (autonomous pipeline)}
\affiliation{NebulaMind Lab, \url{https://lab.nebulamind.net}}
\begin{abstract}
Reionization ionizing-photon-budget reconciliation at z\~{}8: star-forming galaxies require f\_esc=0.299 (+0.301/-0.153) to close the budget (Madau-Dickinson SFRD, log xi\_ion=25.5+/-0.15, clumping C=2-5, JWST-SFRD tail), versus indirect-proxy-inferred f\_esc=0.062 (+0.108/-0.039) from LzLCS O32/beta calibrations. Median delta(required-inferred)=+0.214 dex-frac (16-84\%: +0.042 to +0.517); 90\% of the systematic MC shows a shortfall. Conclusion: a genuine SHORTFALL remains, and the sign holds under both O32 and beta calibrations. The result is bounded by the xi\_ion x clumping x proxy-calibration systematic, not by statistics. Generated autonomously from public data (jwst) via the NebulaMind Lab runner: a bounded, reproducible, descriptive study.
\end{abstract}
\section{Introduction}
Recent studies have highlighted a potential crisis in understanding the reionization process, with concerns that current models may not account for the necessary ionizing photons [Muñoz2024]. This issue is further complicated by the increased demands on ionizing sources due to absorption-dominated reionization [Davies2021]. To address this problem, we revisit the ionizing-photon-budget using a literature-anchored budget calculation.
\section{Data and method}
Our method relies on published values for key parameters: the cosmic star formation rate density (SFRD) is based on the Madau \& Dickinson (2014) analytic fitting function, while xi\_ion and the O32/beta f\_esc proxy calibrations are adopted from recent works [Chisholm+22, Flury+22; Simmonds+24]. We do not utilize any new observational or catalog data, instead focusing on reconciling existing literature values to better understand the reionization process.
\section{Result}
Our analysis reveals a significant shortfall in the ionizing-photon-budget at z\~{}8. Star-forming galaxies would require an escape fraction of f\_esc=0.299 (+0.301/-0.153) to close this budget, assuming the Madau-Dickinson SFRD, log xi\_ion=25.5+/-0.15, and a clumping factor C between 2-5. However, indirect-proxy-inferred escape fractions from LzLCS O32/beta calibrations suggest a much lower value of f\_esc=0.062 (+0.108/-0.039). The median difference between the required and inferred values is +0.214 dex-frac (16-84\%: +0.042 to +0.517), with 90\% of systematic Monte Carlo simulations showing a shortfall.
\begin{figure}\centering\includegraphics[width=\columnwidth]{result.png}\caption{A residual reionization ionizing-photon-budget shortfall for star-forming galaxies at z\~{}8 (lowxi systematic corner)}\end{figure}
\section{Caveats}
It is essential to acknowledge the limitations of our approach, which relies on automated, single-selection, uncalibrated measurements. The accuracy of our results depends heavily on the assumptions and calibrations used in previous studies [Park2022]. Additionally, our analysis does not account for potential uncertainties in the SFRD or xi\_ion values, which could further impact the ionizing-photon-budget. Therefore, while our findings suggest a genuine shortfall, they should be interpreted with caution and considered alongside other observational constraints and model predictions [Madau2017, Duncan2015].
\section*{References}
\footnotesize
[Muoz2024] Muñoz et al. (2024). Reionization after JWST: a photon budget crisis?. bibcode 2024MNRAS.535L..37M \par [Davies2021] Davies et al. (2021). The Predicament of Absorption-dominated Reionization: Increased Demands on Ionizing Sources. bibcode 2021ApJ...918L..35D \par [Park2022] Park et al. (2022). Calibrating excursion set reionization models to approximately conserve ionizing photons. bibcode 2022MNRAS.517..192P \par [Duncan2015] Duncan et al. (2015). Powering reionization: assessing the galaxy ionizing photon budget at z < 10. bibcode 2015MNRAS.451.2030D \par [Madau2017] Madau (2017). Cosmic Reionization after Planck and before JWST: An Analytic Approach. bibcode 2017ApJ...851...50M

\end{document}
